(1) 因 $L=mr^2\dot\phi$，有效势是 $U_{\rm eff}=L^2/(2mr^2)+U(r)$，径向能量积分 $m\dot r^2/2+U_{\rm eff}=E$。对两种势定义都满足 $U'(r)=kr^{p-1}$，所以圆轨道条件
\[
-\frac{L^2}{mr_0^3}+kr_0^{p-1}=0,\qquad L^2=mk r_0^{p+2}.
\]
当 $p\ne-2$，$r_0=(L^2/(mk))^{1/(p+2)}$；当 $p=-2$，只有 $L^2=mk$ 才可能任何半径均为边缘平衡，不能滥用幂式求根。

(2) 二阶导数 $U_{\rm eff}''=3L^2/(mr^4)+k(p-1)r^{p-2}$。在圆轨道代入角动量条件可得
\[
U_{\rm eff}''(r_0)=k(p+2)r_0^{p-2}.
\]
故 $p>-2$ 时稳定，$p<-2$ 不稳定，$p=-2$ 为中性退化、线性频率零（非稳定振动）。此判据对 $p=0$ 的对数势同样有效。

(3) 圆轨道角频率 $\omega_\phi=L/(mr_0^2)=\sqrt{(k/m)r_0^{p-2}}$；小径向振动频率
\[
\omega_r=\sqrt{U_{\rm eff}''(r_0)/m}=\sqrt{p+2}\,\omega_\phi.
\]
近圆极限中相邻近心点的方位角增量为 $\Delta\phi=2\pi/\sqrt{p+2}$。$p=-1$（牛顿 $1/r$ 势）得到 $2\pi$；$p=2$（各向同性谐振势）得到 $\pi$，与经典闭合椭圆轨道相符。一般 $p$ 的近圆轨道不必闭合，$\Delta\phi$ 是小振幅意义而非任意能量轨道的全局结论。

(4) 圆周向心力条件 $mv_\phi^2/r_0=kr_0^{p-1}$ 给 $v_\phi^2=kr_0^p/m$，圆轨道动能 $K=mv_\phi^2/2=kr_0^p/2$。由于 $r^p=1+p\ln r+O(p^2)$，势 $k(r^p-r_*^p)/p\to k\ln(r/r_*)$，导数与频率结果在 $p\to0$ 连续。
